\begin{textsample}{POS Dim 3 – am – Score 21.00 – am/am\_inf\_20.txt}  \label{ex:f3_pos_038}
A Linguagem Escrita Nossa responsabilidade como professor em relação à linguagem escrita na Educação \textbf{Infantil} é de contribuir para a formação de bases para o desenvolvimento dessa linguagem, não é nossa prerrogativa alfabetizar as \textbf{crianças}.
Ao pautarmos nossa prática pedagógica na perspectiva de que devemos contribuir para a formação de bases para o desenvolvimento da linguagem escrita e não na antecipação da alfabetização, é nossa atribuição inicial proporcionarmos às \textbf{crianças} situações e atividades diversas para que elas desenvolvam e exercitem a função simbólica do pensamento, que é uma capacidade essencial para que elas venham a escrever.
A função simbólica do pensamento se caracteriza pela capacidade que a \textbf{criança} desenvolve de usar uma coisa no lugar de outra, fazendo uma substituição, com o objetivo de representar o seu pensamento.
É a partir do desenvolvimento dessa capacidade que se inicia toda a história do desenvolvimento da linguagem escrita da \textbf{criança}, de acordo com Vygotski ( 2012 ).
Segundo Luria ( 2014, p. 143 ), “ a história da escrita na \textbf{criança} começa muito antes que o professor coloque pela primeira vez um \textbf{lápis} em sua mão e lhe ensine a traçar letras ”.
Para o autor, que participou dos estudos de Vygotski, antes que a \textbf{criança} venha a escrever convencionalmente, ela passa por percursos simbólicos que podem ser considerados como uma pré-escrita, que foi denominada como a pré-história da linguagem escrita.
Nesse estudo, Vygotski ( 2012 ) nos aponta que a \textbf{criança}, ao se expressar por meio dos seus \textbf{gestos}, do \textbf{brincar} de faz de conta e do desenho, já está desenvolvendo a sua capacidade de representação, desenvolvendo assim a função simbólica do pensamento.
Em todos esses momentos, a \textbf{criança} se utiliza dessa capacidade como recurso para representar seus pensamentos, o que é básico para que ela venha a escrever posteriormente, já que ao escrevermos, também fazemos esse mesmo processo de representação e de expressão de nossas ideias por meio das palavras.
A pré-história da linguagem escrita se inicia pelo \textbf{gesto} de apontar, quando a \textbf{criança} quer, por exemplo, indicar algo que quer pegar, mas ainda não consegue acessar o objeto de seu \textbf{desejo} sozinho.
O \textbf{gesto} já é uma representação simbólica do seu pensamento, ele representa o \textbf{desejo} ou necessidade da \textbf{criança} de ter algo.
O mesmo também ocorre quando a \textbf{criança} faz seus desenhos e neles procura representar o seu pensamento sobre as coisas.
Ela se utiliza do desenho para representar simbolicamente o que está pensando.
No \textbf{brincar} de faz de conta a \textbf{criança} também faz essa substituição, quando se utiliza de objetos para representar outros objetos ausentes, que usaria numa situação real, para desempenhar seu papel/personagem na \textbf{brincadeira}, como por exemplo, usa um pedacinho de massa de modelar, na ausência de um pedaço de algodão, ao \textbf{brincar}, ou representar um médico.
Em todos esses momentos, a \textbf{criança} está desenvolvendo a função simbólica.
E é nesse sentido, que precisamos como professores e professores da Educação \textbf{Infantil}, ter o conhecimento de que, ao permitirmos que as \textbf{crianças} exercitem diariamente essa capacidade, vivenciando um processo rico de desenvolvimento da gestualidade, da fala, do desenho e do \textbf{brincar}.
Assim, quando ela se aproximar da escrita, já terá os fundamentos necessários de pensamento simbólico para escrever, já que consegue operar com objetos que representam outros objetos e ideias ( os signos ), o que é essencial para escrever ( MORAES, 2015 ).
Concomitantemente, devemos nos dedicar a aproximar as \textbf{crianças} da escrita em sua função social, ou seja, a linguagem escrita deve ser usada a partir de situações reais, para satisfazer a necessidades reais, atendendo, de acordo com Mello ( 2005 ), a finalidade para a qual essa linguagem foi criada : para comunicar, expressar, registrar, informar algo.
A escrita aparece na história da humanidade, inicialmente, como uma linguagem para suprir a necessidade de comunicação entre as pessoas.
Por isso, devemos ter a clareza de que é com essa finalidade que ela deve ser apresentada e usada junto às \textbf{crianças} desde a Educação \textbf{Infantil}, mesmo que as \textbf{crianças} ainda não escrevam \textbf{autonomamente}.
Como podemos fazer isso em nossas \textbf{creches} e \textbf{pré-escolas}?
Primeiramente, estimulando em nossas \textbf{crianças} a necessidade de expressão.
Segundo Mello ( 2007, p.10 ), a necessidade de expressão “ [... ] surge a partir do que as \textbf{crianças} veem, ouvem, vivem, descobrem, e aprendem ”.
Entendemos, assim, que serão as experiências que as \textbf{crianças} terão dentro e fora da escola – na vida real – que demandarão que elas se expressem, inclusive pela linguagem escrita.
Essa expressão pode ser realizada pela fala, pelo \textbf{brincar}, pela \textbf{pintura}, pelo desenho, pela modelagem, pelos movimentos, pelas histórias lidas e \textbf{contadas}, pela dança, pelo cantar, pela produção e construção de objetos, pelos poemas etc., e também pela escrita, quando as \textbf{crianças} externalizam o que pensam sobre o mundo e as coisas por meio da linguagem oral e têm em seu professor ou \textbf{professora} a pessoa que faz a mediação escrevendo a sua frente o que elas falam, assumindo o papel de escriba, organizando essas ideias junto com as \textbf{crianças}, já que elas ainda não dominam a escrita convencional, por isso, ainda não fazem esse processo sozinhas.
Tal atitude é muito bem sintetizada por Britto ( 2005, p. 18 ), quando declara que “ na Educação \textbf{Infantil}, ler com os ouvidos e escrever com a \textbf{boca} ( situação em que a educadora se põe na função de enunciadora ou de escriba ) é mais fundamental do que ler com os olhos e escrever com as próprias mãos ”.
Quando as \textbf{crianças} testemunham esse processo ao verem que seus pensamentos se transformam em linguagem escrita por meio do que o \textbf{adulto} faz para elas estruturando a sua fala em forma de texto, a linguagem escrita está sendo trabalhada de maneira plena diante das \textbf{crianças}.
Nesse processo, elas presenciam a escrita na sua dimensão total, com significado e sentido.
E isso permite que desde o início elas percebam que a escrita é viva e extremamente necessária a vida de todos que habitam a nossa sociedade que se constitui também pela cultura escrita.
Essa postura transforma completamente a visão equivocada pela qual a linguagem escrita vem sendo geralmente tratada na Educação \textbf{Infantil}, por meio de treinos motores da escrita das letras.
A(o) professora/professor precisa compreender que as \textbf{crianças} podem se expressar por escrito com a sua \textbf{ajuda}, ou seja, que elas já produzem textos, fazendo isso por meio das mãos do \textbf{adulto}, não havendo sentido e nem razão para que as \textbf{crianças} dediquem seu tempo a tarefas artificiais de treino motor, já que estarão participando de situações de escrita muito mais completas e complexas que tratam da escrita em sua essência, como linguagem.
Em suma, como nos aponta Jolibert ( 2006, p. 193 ), “ a capacidade de produção de uma \textbf{criança} não deve ser confundida com sua capacidade de grafar suas próprias mensagens ”.
Embora as \textbf{crianças} ainda não grafem com autonomia, já são excelentes produtoras de texto, pois já conseguem expressar o que sentem e pensam com facilidade.
Escrever é, segundo Jolibert ( 2006, p. 191 ), “ produzir mensagens reais, como intencionalidade e destinatários reais.
Não se trata de transcrever ( copiar ) nem de praticar caligrafia [... ] ”, não haverá espaço na Educação \textbf{Infantil} para uma “ escrita ” que não contemple a essência do que seja escrever.
Assim, a prática pedagógica passa por uma transformação profunda no tocante às atividades que serão planejadas e propostas para as \textbf{crianças}, quando se sabe o que significa conceitualmente o que é escrever, considerando que escrever não é copiar letras, e sim expressar o pensamento, as ideias, por meio da linguagem escrita.
Vygotski ( 2012, p. 183 ), em 1931, em seus estudos sobre o desenvolvimento da linguagem escrita, já nos alertava que “ a \textbf{criança} é ensinada a traçar as letras e a formar palavras com elas, mas não se ensina a linguagem escrita ”.
Tal afirmativa nos leva a considerar que o fato de que uma \textbf{criança} consiga preencher várias linhas copiando as letras que lhe são pedidas, não significa que esteja desenvolvendo a linguagem escrita.
A \textbf{cópia} de letras não tem relação com a capacidade de expressão das \textbf{crianças}, afinal, ao se copiar algo estamos reproduzindo o que já está posto, não havendo assim autoria, e a \textbf{criança} não está sendo considerada como sujeito, mas sim uma executora de algo que faz porque alguém lhe manda.
Quando consideramos que as \textbf{crianças} na vida real já convivem, mesmo antes de entrarem na instituição escolar, com a escrita em sua função social, por meio de : mensagens, outdoors, livros, revistas, anúncios, marcas/rótulos de produtos, computadores, celulares, álbuns, jornais, panfletos, comprovantes de pagamento, documentos, encartes de lojas e com muitos outros elementos da cultura escrita, continuaremos seguindo esse mesmo caminho ou perspectiva dentro da instituição, em que a escrita também vai se manifestar \textbf{frequentemente} com função social para satisfazer as necessidades que as \textbf{crianças} apresentarão para usá -la, também para se comunicarem, para se expressarem, para registrarem e informarem algo.
Nesse sentido, Miller e Mello ( 2008 ) nos apontam algumas maneiras pelas quais as \textbf{crianças} podem ir se apropriando da escrita, tendo a/o professora/professor como escriba, como por exemplo : a escrita de cartas para alguém distante ( mas que seja real ) ; a escrita de bilhetes, a confecção de convites ; a elaboração de regras de convivência ; fazer relatórios de passeios ; roteiros de entrevistas ; fazer panfletos de campanhas ; confeccionar cartazes ; registrar descobertas ; fazer o \textbf{registro} da \textbf{rotina} \textbf{diária} etc.
Dando continuidade às atividades propostas pelas autoras, indica -se que as \textbf{crianças} possam fazer uma lista de alimentos que trarão para um lanche coletivo/piquenique ; cartões de felicitações ao se comemorar os aniversariantes da turma ; o \textbf{registro} da frequência das \textbf{crianças} ; o reconto de uma história ; \textbf{registro} de um processo de construção de algum jogo coletivo ; as memórias da turma ; uma receita culinária etc.
Sabemos que muitas outras atividades poderão surgir a partir da realidade e das experiências que cada turma de \textbf{crianças} viverá, não havendo assim uma padronização do que cada uma produz por escrito.
O mais importante é o princípio de que a escrita não se separa da nossa vida real.
Percebemos, a partir dessas atividades, que em todas elas a escrita se apresenta cumprindo a sua finalidade e que todas elas surgem ligadas ao contexto de experiências que as \textbf{crianças} estão tendo na escola, ou seja, a partir de situações reais que emergem das suas vidas no contexto escolar, o que dá significado e sentido para que a linguagem escrita seja utilizada, levandos -as desde muito cedo a perceberem que escrevemos por temos necessidade de escrever.
Nesse sentido, cabe à \textbf{professora} e ao professor mediar intencionalmente as situações para ir mostrando às \textbf{crianças} para que serve a escrita e como ela se apresenta nos diversos materiais e contextos escritos.
Sabemos que muitas das atividades que envolvem a escrita já podem ser previstas pelos professores, porém muitas outras poderão aparecer na dinâmica do dia da dia da escola e não podem ser renegadas por não terem sido planejadas com antecedência.
Devemos sempre lembrar que o que move a prática da escrita junto às \textbf{crianças} é a sua necessidade de uso, por isso, devemos estar atentos ao que acontece para não perdermos ricas oportunidades de aproximar as nossas \textbf{crianças} do universo escrito.

% matched lemmas: adulto, ajuda, autonomamente, boca, brincadeira, brincar, contar, creche, criança, cópia, desejo, diário, frequentemente, gesto, infantil, lápis, pintura, professora, pré-escola, registro, rotina
\end{textsample}
